\section{III. Finite unit certificates without a Neumann smallness hypothesis}\label{iii.-finite-unit-certificates-without-a-neumann-smallness-hypothesis}

\subsection{III.1 The only algebraic inverse gate is the leading coefficient}\label{iii.1-the-only-algebraic-inverse-gate-is-the-leading-coefficient}

\textbf{Source input.} Section 5 of the supplied proof correctly retains the full length-\(m\) physical unit, its derivatives through order \(2m-1\), and the maps between polynomial and physical jets. Its U6--U10 use a sufficient Neumann condition for an uncertain inverse. This section gives a different, finite certificate that uses the triangular jet structure. It does not supply the absent native derivative values.

Let \(S\) be the lower shift in the original divided-jet coordinates of \(\mathbb C[z]/z^m\), so \(S^m=0\). Write the actual unit and its certified centre as \[
 U=\sum_{j=0}^{m-1}u_jS^j,\qquad
 \widehat U=\sum_{j=0}^{m-1}\widehat u_jS^j,\qquad
 |u_j-\widehat u_j|\le e_j.
\] Suppose only that \[
 c:=|\widehat u_0|-e_0>0.
 \tag{U1}
\] A certified positive lower bound may be used in place of the displayed \(c\). Then \(U\) is invertible: it is triangular with diagonal \(u_0\), and its inverse is the \textbf{finite} expression \[
 U^{-1}=u_0^{-1}\sum_{r=0}^{m-1}
 \left(-u_0^{-1}\sum_{j=1}^{m-1}u_jS^j\right)^r.
 \tag{U2}
\] No restriction on the size of the strictly triangular part is needed. Conversely \(u_0=0\) makes \(U\) singular. Thus U1 is the remaining scalar certificate gate, not a suppressed denominator.

Define \(a_j=|\widehat u_j|+e_j\) for \(j\ge1\) and nonnegative majorant coefficients \[
 b_0=c^{-1},\qquad
 b_n=c^{-1}\sum_{j=1}^{n}a_jb_{n-j}\quad(1\le n<m).
 \tag{U3}
\] If \(U^{-1}=\sum v_nS^n\), comparison of coefficients of \(UU^{-1}=I\) proves inductively \(|v_n|\le b_n\). Compute the centre inverse \(\widehat v_n\) by the same exact signed convolution, not by a numerical matrix inverse. From \[
 U^{-1}-\widehat U^{-1}=-U^{-1}(U-\widehat U)\widehat U^{-1}
\] we obtain coefficient radii \[
 \boxed{|v_n-\widehat v_n|\le f_n:=\sum_{i+j+l=n}b_i e_j|\widehat v_l|.}
 \tag{U4}
\] All sums stop before degree \(m\).

\subsection{III.2 The retained jet metric improves rather than destroys the estimate}\label{iii.2-the-retained-jet-metric-improves-rather-than-destroys-the-estimate}

Let \(D=\operatorname{diag}(d_0,\ldots,d_{m-1})>0\) be the specified polynomial-jet Gram. Define \[
 \beta_j=\|S^j\|_D
 =\max_{0\le l<m-j}\sqrt{d_{l+j}/d_l},\qquad \beta_0=1.
 \tag{U5}
\] The formula follows because \(S^j\) sends each coordinate to a different coordinate; its weighted squared norm is the maximum of the displayed ratios. Therefore \[
 \|U^{-1}\|_D\le\sum b_j\beta_j,\qquad
 \|U^{-1}-\widehat U^{-1}\|_D\le\sum f_j\beta_j.
 \tag{U6}
\] These are finite bounds even when an unweighted Neumann test fails.

For the supplied binomial jet family, \(d=m-1\) and \[
 d_l=a(l!)^2\binom{m-1}{l}\eta^{2l},\qquad a>0,\quad\eta>0.
\] The exact ratio is \[
 \boxed{\beta_j=\eta^j\max_l
 \sqrt{\frac{(l+j)!}{l!}\frac{(m-1-l)!}{(m-1-l-j)!}}
 \le(\eta m/2)^j.}
 \tag{U7}
\] Indeed every one-step factor is \(\eta\sqrt{(l+1)(m-1-l)}\le\eta m/2\), and \(S^j\) is a product of \(j\) shifts. In particular higher unit coefficients are suppressed by their actual powers of \(\eta\) in this metric. This statement concerns the specified constructed jet metric, not the unrelated canonical scalar-source metric.

\subsection{III.3 Two-sided physical-metric bounds with no smallness restriction}\label{iii.3-two-sided-physical-metric-bounds-with-no-smallness-restriction}

Retain \[
 H=U^{-*}DU^{-1},\qquad
 \widehat H=\widehat U^{-*}D\widehat U^{-1}.
\] Put \[
 \zeta=\sum_{n=0}^{m-1}\beta_n\sum_{i+j=n}b_i e_j,
 \qquad
 \rho=\sum_{n=0}^{m-1}\beta_n\sum_{i+j=n}|\widehat v_i|e_j.
 \tag{U8}
\] In the original physical coordinates the complete bound is \[
 \boxed{(1+\rho)^{-2}\widehat H
 \preceq H\preceq(1+\zeta)^2\widehat H.}
 \tag{U9}
\] \textbf{Proof.} Let \(R=U^{-1}\widehat U\). Then \[
 R=I-U^{-1}(U-\widehat U),\qquad
 R^{-1}=I+\widehat U^{-1}(U-\widehat U).
\] Finite convolution and U5 show \(\|R\|_D\le1+\zeta\) and \(\|R^{-1}\|_D\le1+\rho\). Thus \[
 (1+\rho)^{-1}\|y\|_D\le\|Ry\|_D\le(1+\zeta)\|y\|_D.
\] Take \(y=\widehat U^{-1}x\) and square. This proves U9 without requiring \(\rho<1\) or \(\zeta<1\). If \(\zeta<1\), the lower bound may be sharpened to the larger of \((1-\zeta)^2\) and \((1+\rho)^{-2}\). If \(\rho<1\), the upper bound may be sharpened to the smaller of \((1+\zeta)^2\) and \((1-\rho)^{-2}\).

\textbf{A finite example outside the old sufficient domain.} Let \(m=3\), \(\widehat U=I\), \(U=I+2S\), \(e_1=2\), \(e_0=e_2=0\), and \(\eta=1/100\). The unweighted error test \(\|\widehat U^{-1}\|\sum e_j<1\) fails because its left side is two. The actual inverse is \(I-2S+4S^2\). With \[
 D=\operatorname{diag}(1,2\eta^2,4\eta^4),\qquad b=3\eta/2,
\] U7--U8 permit the rational bounds \[
 \rho=2b=3/100,\qquad
 \zeta=2b+4b^2=309/10000.
\] Hence \[
 (100/103)^2D\preceq U^{-*}DU^{-1}
 \preceq(10309/10000)^2D.
 \tag{U10}
\] The inverse exists with the displayed close relative metric bounds even though the generic Neumann gate failed. This is a declared triangular diagnostic, not a computation of a native unit.

\subsection{III.4 Tensoring, quotient minima, and the original labels}\label{iii.4-tensoring-quotient-minima-and-the-original-labels}

For factorwise bounds U9, positivity of tensor products gives \[
 \boxed{\left(\prod_i(1+\rho_i)^{-2}\right)
 \bigotimes_i\widehat H_i
 \preceq\bigotimes_iH_i
 \preceq
 \left(\prod_i(1+\zeta_i)^2\right)
 \bigotimes_i\widehat H_i.}
 \tag{U11}
\] Restriction by the \textbf{same} cyclic injection and minimization over the \textbf{same} affine value fibres preserve both scalar factors: compare the squared norm of every admissible representative and then take its infimum. A \(q\)-dimensional attained determinant therefore has absolute logarithmic error at most \[
 \boxed{2q\max\left\{\sum_i\log(1+\rho_i),
                        \sum_i\log(1+\zeta_i)\right\}.}
 \tag{U12}
\] When approximate units are used as value labels, the value map, observation and class must first be transported by the actual tensor product of \(\widehat U_i^{-1}U_i\). Equation U11 is not permission to change those fibres silently.

For a fixed factor with an exact nonzero leading unit and bounded higher coefficient uncertainties, U7 implies \(\rho,\zeta=O(\eta)\) as \(\eta\downarrow0\). Tensoring \(k\) such factors gives a displayed exponential factor \(\exp(O(k\eta))\), not an assumed equality of the metrics. With uncertain leading coefficient there is also its nonvanishing scalar contribution; it does not acquire an \(\eta\) factor.

The input's derivative conversion remains indispensable: for a root of complete multiplicity \(m\), the unit jet uses \(g^{(m)},\ldots,g^{(2m-1)}\), and the full local cofactor. The present certificate changes only the inverse-error propagation after those primitive derivative enclosures have actually been supplied.
